\chapter{Composite Systems and Local Tomography}
\label{ch:local-tomography}

Part I's two axes — compositional sufficiency (Chapter~\ref{ch:composition})
and global realizability (Chapters~\ref{ch:context-graphs}--\ref{ch:real-mub-obstruction})
— both concern a \emph{single} system observed through different contexts or
composed sequentially with itself. This chapter opens the third axis: given
two \emph{separate} systems, individually admissible in the sense of Part I,
how must their joint description compose?

\section{The Composition Problem}

Let systems $A,B$ have individual state spaces of dimension $d_A,d_B$
respectively (each already subject to the admissibility constraints of
Chapters~\ref{ch:simplex}--\ref{ch:unistochastic}). Operational independence
means: $A$ and $B$ can be prepared and measured separately, and a joint
preparation of both is a legitimate object of study without presupposing in
advance what mathematical structure describes it. The question this chapter
asks is deliberately neutral:
\begin{equation}
\boxed{\text{How must independently accessible systems compose?}}
\end{equation}
rather than assuming a complex tensor product from the outset and asking only
whether it is consistent.

\section{\texorpdfstring{$\mathbb F$}{F}-Hermitian State Spaces}

For each candidate scalar field $\mathbb F \in \{\R,\C,\Hset\}$, a natural
(not yet composite) state space for a $d$-dimensional system is the space of
$\mathbb F$-Hermitian $d\times d$ matrices $\rho$ (density operators
$\rho\ge0$, $\operatorname{tr}\rho=1$ living inside this linear space), where
$\mathbb F$-Hermitian means $\rho^\dagger=\rho$ with the conjugate taken in
$\mathbb F$.

\begin{proposition}[Dimension of the Hermitian state space]
\label{prop:hermitian-dimension}
\index{Dimension of the Hermitian state space (proposition)}
The real vector space of $\mathbb F$-Hermitian $d\times d$ matrices has
dimension
\begin{equation}
\beta_\R(d) = \frac{d(d+1)}{2}, \qquad
\beta_\C(d) = d^2, \qquad
\beta_\Hset(d) = d(2d-1).
\end{equation}
\end{proposition}

\begin{proof}
In each case, Hermiticity forces the diagonal entries to be real (a
self-conjugate element of $\R,\C,\Hset$ is real), contributing $d$ real
parameters, while each of the $d(d-1)/2$ upper-triangular off-diagonal entries
is an unconstrained element of $\mathbb F$ (the lower triangle being
determined by conjugating the upper), contributing $\dim_\R(\mathbb F)$ real
parameters each: $1$ for $\R$, $2$ for $\C$, $4$ for $\Hset$. Summing,
\begin{equation}
\beta_\R(d) = d + 1\cdot\tfrac{d(d-1)}{2} = \tfrac{d(d+1)}{2}, \qquad
\beta_\C(d) = d + 2\cdot\tfrac{d(d-1)}{2} = d^2, \qquad
\beta_\Hset(d) = d + 4\cdot\tfrac{d(d-1)}{2} = d(2d-1).
\end{equation}
\end{proof}

\section{Local Tomography}

\begin{definition}[Local tomography, dimension-counting form]
A scalar field $\mathbb F$ satisfies \emph{local tomography} if, for every
$d_A,d_B$, the Hermitian state space of the composite system has dimension
equal to the product of the local dimensions:
\begin{equation}
\beta_{\mathbb F}(d_A d_B) = \beta_{\mathbb F}(d_A)\, \beta_{\mathbb F}(d_B).
\end{equation}
Informally: the joint state is exactly as informative as a tomographically
complete set of \emph{local} measurements on each subsystem, with no
additional holistic parameters and no missing ones.
\end{definition}

\begin{theorem}[Local tomography selects $\C$]
\label{thm:local-tomography}
\index{Local tomography selects C (theorem)}
For all $d_A,d_B \ge 2$:
\begin{equation}
\beta_\C(d_Ad_B) - \beta_\C(d_A)\beta_\C(d_B) = 0,
\end{equation}
\begin{equation}
\beta_\R(d_Ad_B) - \beta_\R(d_A)\beta_\R(d_B) = \frac{d_Ad_B(d_A-1)(d_B-1)}{4} > 0,
\end{equation}
\begin{equation}
\beta_\Hset(d_Ad_B) - \beta_\Hset(d_A)\beta_\Hset(d_B) = -2\,d_Ad_B(d_A-1)(d_B-1) < 0.
\end{equation}
\end{theorem}

\begin{proof}
Direct substitution of Proposition~\ref{prop:hermitian-dimension} into each
difference, using $d_Ad_B - d_A - d_B + 1 = (d_A-1)(d_B-1)$:
\begin{align}
\beta_\C(d_Ad_B) - \beta_\C(d_A)\beta_\C(d_B) &= (d_Ad_B)^2 - d_A^2 d_B^2 = 0,\\
\beta_\R(d_Ad_B) - \beta_\R(d_A)\beta_\R(d_B) &= \frac{d_Ad_B(d_Ad_B+1)}{2} -
\frac{d_A(d_A+1)}{2}\cdot\frac{d_B(d_B+1)}{2} = \frac{d_Ad_B(d_A-1)(d_B-1)}{4},\\
\beta_\Hset(d_Ad_B) - \beta_\Hset(d_A)\beta_\Hset(d_B) &= d_Ad_B(2d_Ad_B-1) -
d_A(2d_A-1)d_B(2d_B-1) = -2d_Ad_B(d_A-1)(d_B-1).
\end{align}
For $d_A,d_B \ge 2$, $(d_A-1)(d_B-1) > 0$, giving the stated signs.
\end{proof}

\begin{remark}
The mismatch is not a boundary artifact of small dimensions: it grows without
bound as $d_A,d_B$ increase, in opposite directions for $\R$ and $\Hset$. Real
composite systems always carry \emph{more} joint parameters than local
tomography predicts (extra, non-locally-accessible correlational structure);
quaternionic composite systems always carry \emph{fewer} (the naive product
construction is already overdetermined, consistent with the well-known
difficulty of defining a well-behaved tensor product for quaternionic Hilbert
spaces at all). Complex Hermitian state spaces are the unique case among the
three where the identity holds exactly, for every pair of dimensions.
\end{remark}

\begin{scopebox}
\scopetitle{}"Unique" above means unique \emph{among $\R,\C,\Hset$, under
this specific dimension-counting test, using the ordinary Kronecker
product}. It is not a claim that $\C$ is the only conceivable scalar
structure for physics, and it is not unconditional: local tomography
itself is an \emph{assumed} principle this chapter tests rather than
derives (this chapter's closing ledger says so explicitly), and
Chapter~\ref{ch:final-audit}'s \S\ref{sec:recent-literature} discusses
published work in which changing the composition rule (not the scalar
field) lets a real theory pass an equivalent test. Read this theorem as
"\emph{given} local tomography and the Kronecker product, $\C$ is the
unique fixed point" — every word of that qualification is load-bearing.
\end{scopebox}

\begin{example}[$d_A=d_B=2$]
Two real "qubits": $\beta_\R(4)=10$ vs.\ $\beta_\R(2)^2=9$ — one real parameter
of the joint two-qubit state is invisible to any local tomography scheme.
Two quaternionic "qubits": $\beta_\Hset(4)=28$ vs.\ $\beta_\Hset(2)^2=36$ — the
naive quaternionic product already overcounts by eight real parameters before
any physical constraint is imposed. Two ordinary qubits: $\beta_\C(4)=16 =
\beta_\C(2)^2$ exactly.
\end{example}

\begin{ledgerentry}[Scope of this chapter]
\textbf{Established:} local tomography, precisely defined as a
dimension-counting identity on Hermitian state spaces, holds identically for
$\C$ and fails identically (in opposite directions) for $\R$ and $\Hset$, for
every pair of dimensions $d_A,d_B \ge 2$
(Theorem~\ref{thm:local-tomography}). This argument is not original to this
book; it is the classical parameter-counting characterization due to
Araki~\cite{araki1980}, sharpened by Wootters and later incorporated into
Hardy's axiomatic reconstruction of quantum theory~\cite{hardy2001}.\\
\textbf{Not established:} that local tomography is itself \emph{forced} by
anything proved in Part I — it is introduced here as an additional principle
under test, in keeping with the book's assumption-ledger methodology, not as
a consequence of the single-system results. Nor does dimension-counting alone
establish the full tensor-product structure, the Born rule for composite
measurements, or entanglement; it constrains the \emph{dimension} of the
composite state space and no more. What Part I's cycle obstruction
(Chapter~\ref{ch:real-mub-obstruction}) and this chapter's local-tomography
argument have in common — ruling out $\R$ — arise by entirely different
mechanisms (contextual gluing around a cycle vs.\ parameter counting across a
tensor product) and should be credited as independent lines of evidence, not
conflated into one argument.
\end{ledgerentry}
